A father is on the road home. When he left his house, his son was dying.
He went to find Jesus because his child needed help that he could not
give. Now he is returning, but Jesus is not walking beside him.
The father has a word to carry home.

That is the moment to stay with in today's Gospel. Before the servants
arrive, before the father learns when the fever broke, he begins the
journey. What has changed? He has not seen his child. He has heard
Jesus, and the word he has heard is enough to set him moving.

At first, the father asks Jesus to come down and heal his son.
There is urgency in the request. The child is at the point of death.
When Jesus challenges a faith dependent on signs and wonders, the
father asks again: ``Lord, come down before that my son die.''
He wants Jesus in the room. He wants the healer where the sickness is.

Saint Gregory the Great notices something here. The father must believe
that Jesus can help, or he would not have come. Yet he still thinks
that Christ's power needs Christ's bodily presence beside the child.
His trust is real, but his understanding of the one he trusts is too
small. The Lord who made the world does not lose his power across the
distance between two towns.

Jesus answers: ``Go thy way. Thy son liveth.'' The father came asking
Jesus to make a journey. Jesus gives the father a journey to make.
And John tells us that the man believed the word and went his way.
Christ heals the son, and he teaches the father to trust him.
The gift and the lesson arrive together.

Then the servants meet him. They bring good news. The father asks
when the child began to recover, and the hour they name is the hour
when Jesus spoke. John wants us to hear that correspondence. The
healing really happened; the testimony matters. Faith does not make
the father's question unnecessary. But his departure came before
the answer. He had already begun to act on Christ's word.

We can postpone obedience until we feel we possess the whole outcome.
Think of an admission of fault that we know we owe someone. We may
wait until we can be certain the conversation will end well. The
other person must first become receptive; our own embarrassment must
first disappear. Meanwhile the needed words go unspoken. Our desire
to control the answer has begun to govern whether we will tell the truth.

Today's Epistle brings this movement into our daily lives. It tells
us to walk wisely, to understand God's will, and to redeem the time
because the days are evil. Wisdom is needed precisely when the
circumstances are difficult. We do not have to wait for a trouble-free
world before we can be faithful within it.

Saint John Chrysostom explains that such wisdom may cost us an
advantage. We may have to give up the satisfaction of answering
hostility with hostility. Keeping faith can matter more than winning
the encounter. An honest admission may leave us with less to defend and
more to repair. Yet we can begin that work because truthfulness
belongs to the Lord's will, even while another person's response
remains beyond our control.

The Epistle moves through song and thanksgiving and ends with being
subject to one another in the fear of Christ. It reaches the way we
treat people. The heart that sings to God is learning to serve the
neighbor. Belief in Christ's word has something to do with what we
say at home and whose burden we help to carry.

But there is a hard question beside the road home. What about the
father or mother who prays and does not receive this good news?
What about the sick person who remains sick? This father received
an explicit assurance from Jesus that his son lived. We have not
been given that same assurance about every illness. This particular
cure is not a promise that sufficient faith will secure whatever
bodily recovery we ask for.

The Mass itself keeps us from making that mistake. At the Offertory
it gives us the captives beside the rivers of Babylon, remembering
Zion and weeping. Their tears belong in the worship of God. At
Communion it gives us the servant who finds comfort in God's promise
while still in humiliation. The suffering has not first disappeared.
Hope can remain in a life that still hurts.

Saint Augustine explains that when this servant asks God to remember
his word, he is asking for the promise to be fulfilled. God has not
forgotten him. The promise gives him something to hold to within
his trial. Augustine carries that hope as far as the Church's
resurrection. Christ's gift reaches beyond the span of a recovered
earthly life to life with God. That hope leaves us free to seek help,
to care for the sick, and to grieve with those who grieve. None of
those acts requires us to predict the outcome.

Now listen again to the father asking the servants about the hour.
Their answer changes his understanding of the journey he has just
made. While he was still on the road, unable to see his son, the
child was already well. The healing had not waited for his arrival.
What he learns now is what Christ had done when he spoke.

That is a greater discovery than finding that he had managed to be
brave. His confidence has an object: the Lord who gave life to his
son. The matching hour directs us to the power and mercy of the
one he believed.
He has reason to give thanks. He can name the gift and the giver.

This matters when we hear the Church ask for healing at the altar.
The Secret asks that these mysteries give heavenly medicine and
cleanse the vices of our hearts. We need more than an inspiring
example of someone who acted courageously. We need Christ himself
to heal what makes our love crooked. The Eucharistic gift answers
that need.

The Collect asks for pardon and peace so that we may serve God with
a secure mind. That peace has its source in received mercy. When
we acknowledge a fault, we do not have to protect an imagined
innocence at all costs. We have a Savior from whom to seek pardon.
And when the Church asks after Communion that we may obey God's
commandments, she asks for help to live from his gift. The father
was sent home by the Lord's word; the faithful are sustained in
obedience by the Lord whose gifts they receive.

John's last sentence opens the door of the whole house. The father
believes, and his household believes. The healing has brought more
than relief to one anxious man. Christ has become the ground of
a shared faith. The father sought help for his son, and the mercy
has reached people beyond the two of them.

Notice how the servants belong to that ending. The father receives
their testimony. These people who serve in his house know something
he needs to hear. He asks them, listens to them, and recognizes
Christ's action through what they tell him. Faith brings him into
a deeper dependence on the Lord, and the story also gives him a
reason to attend to the people around him. The Epistle's mutual
service has a home to enter.

We cannot promise our households this father's cure. We can bring
home the truth on which his thanksgiving rests: mercy comes from
Christ, and it can become a good that others share. An honest
admission, freed from the need to defend ourselves, can let the
people closest to us encounter the fruit of pardon. The faith
we profess begins to change the life we have together.

The father went home believing a word. He learned on the road that
Christ's mercy had already reached his child. At the end, the house
itself is filled with faith. Let our homes receive the fruit of
what Christ gives us here. Let those who live with us find people
who know their need for mercy, and who have learned where mercy
is found.
